\documentclass[a4paper, 12pt]{bookln9}
\usepackage[utf8]{inputenc}
\usepackage[spanish]{babel}
\decimalpoint % To use point as decimal point
\usepackage{enumerate}
\usepackage{graphicx}
\usepackage{hyperref}
\usepackage{cite}
\usepackage{textcomp}
\usepackage{subfig}
\usepackage{multirow}
\usepackage{ntheorem}
\usepackage{hyperref}
\usepackage[dvipsnames]{xcolor}
\usepackage{color}
\usepackage{float}
\hypersetup{
	colorlinks = true,
	linkcolor = darkblue,
	anchorcolor = darkblue,
	citecolor = darkblue,
	filecolor = darkblue,
	urlcolor = darkblue
}
\usepackage{amsmath}
\usepackage{amssymb}
\usepackage{lscape}
\usepackage{url}

 % Para cajas con código R
\usepackage{listings}
\lstset{basicstyle=\ttfamily, breaklines=true}
\lstset{ 
	language=R,                     % the language of the code
	basicstyle=\footnotesize\ttfamily, % the size of the fonts that are used for the code
	numbers=left,                   % where to put the line-numbers
	numberstyle=\footnotesize\color{gray},  % the style that is used for the line-numbers
	stepnumber=1,                   % the step between two line-numbers. If it is 1, each line
	% will be numbered
	numbersep=5pt,                  % how far the line-numbers are from the code
	backgroundcolor=\color{white},  % choose the background color. You must add \usepackage{color}
	showspaces=false,               % show spaces adding particular underscores
	showstringspaces=false,         % underline spaces within strings
	showtabs=false,                 % show tabs within strings adding particular underscores
	frame=single,                   % adds a frame around the code
	rulecolor=\color{black},        % if not set, the frame-color may be changed on line-breaks within not-black text (e.g. commens (green here))
	tabsize=2,                      % sets default tabsize to 2 spaces
	captionpos=b,                   % sets the caption-position to bottom
	breaklines=true,                % sets automatic line breaking
	breakatwhitespace=false,        % sets if automatic breaks should only happen at whitespace
	keywordstyle=\color{RoyalBlue},     % keyword style
	commentstyle=\color{Mulberry},      % comment style
	stringstyle=\color{ForestGreen},    % string literal style
	literate={á}{{\'a}}1
	{ã}{{\~a}}1
	{é}{{\'e}}1
	{ó}{{\'o}}1
	{í}{{\'i}}1
	{ñ}{{\~n}}1
	{Ñ}{{\~N}}1
	{¡}{{!`}}1
	{¿}{{?`}}1
	{ú}{{\'u}}1
	{Í}{{\'I}}1
	{Ó}{{\'O}}1
} 


\lstset{
	language = VBScript,
	literate={á}{{\'a}}1
	{ã}{{\~a}}1
	{é}{{\'e}}1
	{ó}{{\'o}}1
	{í}{{\'i}}1
	{ñ}{{\~n}}1
	{Ñ}{{\~N}}1
	{¡}{{!`}}1
	{¿}{{?`}}1
	{ú}{{\'u}}1
	{Í}{{\'I}}1
	{Ó}{{\'O}}1
} 


\usepackage{fancy} % cls file
\usepackage{mathfont} % cls file

\definecolor{darkblue}{rgb}{0.0, 0.0, 0.55}



\newtheorem*{defi}{\normalfont\fontfamily{phv}\fontsize{12}{17}\bfseries Definici\'on}[section]
\newtheorem*{lema}{\normalfont\fontfamily{phv}\fontsize{12}{17}\bfseries Lema}[section]
\newtheorem*{conj}{\normalfont\fontfamily{phv}\fontsize{12}{17}\bfseries Conjetura}[section]
\newtheorem*{coro}{\normalfont\fontfamily{phv}\fontsize{12}{17}\bfseries Corolario}[section]
\newtheorem*{car}{\normalfont\fontfamily{phv}\fontsize{12}{17}\bfseries Caracterización}[section]
\newtheorem{propos}{\normalfont\fontfamily{phv}\fontsize{12}{17}\bfseries Proposición}[section]
\newtheorem*{prop}{\normalfont\fontfamily{phv}\fontsize{12}{17}\bfseries Propiedades}[section]
\newtheorem{nota}{\normalfont\fontfamily{phv}\fontsize{12}{17}\bfseries Nota}[section]
\newtheorem*{ejemplo}{\normalfont\fontfamily{phv}\fontsize{12}{17}\bfseries Ejemplo}[section]
\newtheorem{teorema}{\normalfont\fontfamily{phv}\fontsize{12}{17}\bfseries Teorema}[section]
\newtheorem*{demostracion}{\sc Demostración} 

\renewcommand{\thefootnote}{\arabic{footnote}}
%\renewcommand{\labelitemi}{$\cdot$}
\renewcommand\labelitemi{-} % - en itemize
%\renewcommand{\ttdefault}{phv}

%%%%%%%%%%%%%%%%%%%%%%% MEDIDAS %%%%%%%%%%%%%%%%%%%%%%%%%%%%%%%

% \parindent      0 mm
 \topmargin      10 mm   %%%%%% 0mm
 \headsep        8 mm   %%%%%% 8mm
 \headheight     5 mm   %%%%%% 5mm
 \textheight     211 mm
 \footskip       10 mm
 \headrulewidth  0.7pt
 \oddsidemargin  5 mm
 \evensidemargin  29.4 mm
\evensidemargin 5 mm
\textwidth      145 mm
\renewcommand{\baselinestretch}{1.15}
%%%%%%%%%%%%%%%%%%%%%%%%%%%


\setlength\parindent{0pt} %% Eliminar indent (sangría) en todo el documento

\begin{document}

%%%%%%%%%%%%%%%%%% HEADINGS %%%%%%%%%%%%%%%%%%%%%%%%%%%%%%%%%%%
 \pagestyle{fancyplain}
 \lhead[\fancyplain{}{\small\bf\thepage}]{\fancyplain{}{\small\bf\rightmark}}
 \rhead[\fancyplain{}{\small\bf\leftmark}]{\fancyplain{}{\small\bf\thepage}}
 \cfoot[\fancyplain{\small\bf\thepage}{}]{\fancyplain{\small\bf\thepage}{}}
%%%%%%%%%%%%%%%%%%%%%%%%%%%%%%%%%%%%%%%%%%%%%%%%%%%%%%%%%%%%%%%	

\author{\textbf{Autor:}\vspace{-0.1cm}\\ Daniel Redondo Sánchez \vspace{0.2cm}\\
\textbf{Tutores:}\vspace{0.1cm}\\
Ignacio Rojas\\
Luis Javier Herrera\\
Daniel Castillo
% \newline \textsc{Universidad de Granada}
}
	
\title{
	\vspace{-4cm}
	\centering
	{\Large TRABAJO FIN DE MÁSTER\\} \vspace{0.3cm}
	{\large \textsc{Máster Universitario Oficial en Ciencia de Datos e Ingeniería de Computadores\\}}
		{\LARGE \textbf{\bfseries{Epidemiología y detección de biomarcadores en cáncer\\}}}
		\vspace{-0.75cm}
}


\date{\vspace{1cm}
	Granada, septiembre de 2020 \\
	\vspace{0.5cm}
	\includegraphics[height=2.5cm]{logos/ugr.png} \\ 
}
	
	\mainmatter
	\maketitle
	\thispagestyle{empty}
	
	%% 
	%% \null\newpage
	%% \thispagestyle{empty}
	%% \noindent \textbf{Agradecimientos:}\\
	%% 
	%% 
	%% \null\newpage
	%%
	%% 
	
	% Índice
	\tableofcontents
		
	% Abstract
	\thispagestyle{plain}

\markboth{Abstract}{Abstract}

\vspace{-30pt}

\section*{Resumen}
\addcontentsline{toc}{chapter}{Resumen / Abstract}

\textbf{Introducción:} El cáncer es uno de los mayores problemas de salud pública del mundo con más de 17 millones de casos nuevos y 9 millones de defunciones al año. \\

\textbf{Métodos:} Este trabajo se centra en el cáncer de hígado y el cáncer de colon-recto, describiendo sus principales indicadores epidemiológicos y usando \textit{machine learning} para analizar más de 1.100 muestras de RNA-Seq procedentes de pacientes de cáncer. Para clasificación biclase (tumor vs. tejido normal) y multiclase (varios tipos de tumor vs. tejido normal) se identifican los 10 genes más relevantes y se construyen modelos predictivos con SVM, \textit{random forest} y kNN con validación cruzada 5-fold.\\

\textbf{Resultados:} Los mejores clasificadores biclase para cada algoritmo son validados con excelentes resultados para cáncer de hígado (F1-Score en test: 99,5\%) y cáncer de colon-recto (F1-Score: 100\%). En los mejores modelos para clasificación multiclase se obtienen medidas de evaluación inferiores tanto en hígado (F1-Score: 91,8\%) como en colon-recto (F1-Score: 79,3\%).\\

Se ha desarrollado una aplicación web, \texttt{\textit{biomarkeRs}}, que implementa análisis de transcriptómica y puede resultar de utilidad para usuarios sin conocimientos previos de programación.\\

\textbf{Conclusiones:} SVM, random forest y kNN obtienen resultados muy similares, y consiguen distinguir correctamente entre tejidos tumorales y sanos, con algunos problemas para distinguir entre diferentes tipos de cáncer. Es necesaria una validación externa e interpretaciones clínicas para establecer de forma clara una asociación gen-enfermedad.\\

\textbf{Palabras clave:} epidemiología, transcriptómica, cáncer de hígado, cáncer de colon-recto, RNA-Seq, machine learning, SVM, random forest, kNN.

\newpage
\thispagestyle{plain}

\section*{Abstract}

\textbf{Introduction:} Cancer is one of the world's largest public health problems with more than 17 million new cases and 9 million deaths every year. \\

\textbf{Methods:} This work focuses on liver cancer and colon-rectum cancer, describing their main epidemiological indicators and using machine learning to analyze more than 1,100 RNA-Seq samples from cancer patients. For binary (tumor vs. normal tissue) and multiclass (various tumor types vs. normal tissue) classification, the 10 most relevant genes are identified and predictive models are constructed with SVM, random forest and kNN with 5-fold cross-validation. \\

\textbf{Results:} The best binary classifiers are validated with excellent results for liver cancer (F1-Score in test: 99.5\%) and colon-rectum cancer (F1-Score: 100\%). Lower evaluation measures are obtained in the best models for multiclass classification, in both liver (F1-Score: 91.8\%) and colon-rectum (F1-Score: 79.3\%). \\

A web application has been developed, \texttt{\textit{biomarkeRs}}, that implements transcriptomic analysis and can be useful for users with no previous knowledge of programming. \\

\textbf {Conclusions:} SVM, random forest and kNN obtained very similar results, and managed to correctly distinguish between tumoral and normal tissues with some troubles distinguishing between different types of cancer. External validation and clinical interpretations are necessary to clearly establish a gene-disease association.\\

\textbf{Keywords:} epidemiology, transcriptomics, liver cancer, colorectal cancer, RNA-Seq, machine learning, SVM, random forest, kNN.

\newpage
\thispagestyle{empty}


	
	% Capítulos
	
	% Introducción
	\chapter{Introducción}

\section{Cáncer}


\section{Cáncer de páncreas}


\section{Causas del cáncer de páncreas}

\textcolor{red}{[Prevención]}

	
	% Epidemiología del cáncer de páncreas
	\chapter{Epidemiología del cáncer}

Definición epidemiología \cite{IARC1999}.

\section{Indicadores epidemiológicos}

Para medir en la población el impacto del cáncer se utilizan principalmente cuatro indicadores:

\begin{itemize}
	\item \textbf{Incidencia} (casos nuevos). Mide el riesgo de presentar cáncer.
	\item \textbf{Mortalidad} (defunciones). Mide el riesgo de morir por cáncer.
	\item \textbf{Supervivencia} (porcentaje de casos vivos). Mide la historia natural del cáncer y efectividad del tratamiento.
	\item \textbf{Prevalencia} (casos nuevos y antiguos, vivos). Mide la carga asistencial de la enfermedad.
\end{itemize}

\textcolor{red}{AÑADIR TENDENCIAS!! DE TODO}

\section{Epidemiología del total del cáncer (¿excepto piel no melanoma?)}

\textcolor{red}{Medidas de incidencia}

\textcolor{red}{Cómo se obtiene la incidencia. Importancia de RCPoblacionales, estimaciones y proyecciones. Para incidencia, la unidad de análisis es del tumor, no la persona.}

\textcolor{red}{Número de casos. Tiene problemas por tamaño de población. Tasa bruta. Tiene problemas por estructura de población. Tasas estandarizadas (mundiales, europeas viejas y nuevas). Se usan a veces otros indicadores como tasas acumulativas.}


\textcolor{red}{Diagrama de Marimekko de incidencia de cáncer. Añadir categoría de Otros}

\textcolor{red}{Cómo se obtiene la mortalidad. Importancia de certif de defunción. También estimaciones y proyecciones.}

\textcolor{red}{Supervivencia se calcula principalmente a partir de inc, mort y tablas de vida población general}

\section{Epidemiología del cáncer de hígado}

\subsection{Incidencia de cáncer de hígado}

\subsection{Mortalidad por cáncer de hígado}

\subsection{Supervivencia de cáncer de hígado} 

\subsection{Prevalencia de cáncer de hígado}

\section{Epidemiología del cáncer de colon-recto}

\subsection{Incidencia de cáncer de hígado}

\subsection{Mortalidad por cáncer de hígado}

\subsection{Supervivencia de cáncer de hígado} 

\subsection{Prevalencia de cáncer de hígado}







	
	% Epidemiología del cáncer de páncreas
	\chapter{\textit{Machine learning} aplicado a transcriptómica}

\section{Selección de biomarcadores}

\textcolor{red}{mRMR, RF, DA}

\section{Algoritmos de clasificación}

\textcolor{red}{SVM; KNN}

	
	% Machine learning aplicado a genómica
    \chapter{Detección de biomarcadores en cáncer de páncreas}

\section{Introducción}

\section{Metodología}

\textcolor{red}{Fuente de los datos, etc}

\section{Resultados}


\section{Conclusiones}


	
	% Detección de biomarcadores en cáncer de páncreas
	\chapter{Detección de biomarcadores en cáncer de páncreas}

\section{Introducción}

\section{Metodología}

\textcolor{red}{Fuente de los datos, etc}

\section{Resultados}


\section{Conclusiones}


	
	% Detección de biomarcadores en cáncer de colon-recto
	\chapter{Aplicación web para realizar detección de biomarcadores}

	
	% Conclusiones y líneas abiertas de trabajo
	\chapter{Líneas abiertas de trabajo}




	
	% Bibliografía
	\newpage
	\phantomsection
	\addcontentsline{toc}{chapter}{Bibliografía}
	\bibliographystyle{unsrt}
	\thispagestyle{empty}
	\bibliography{bibliografia}
	
	% Apéndices
	\appendix
	\backmatter
	\chapter{AnexoI: Código de análisis en R}
	\section*{Anexo: Código en R}\label{anexo}

\subsection*{Funciones}

\textbf{Caja XX}. Definición de \texttt{funcion}.

\lstinputlisting[language=R]{codigo/cabecera.R}

	\chapter{Anexo II: Código de aplicación web}
	\section*{Anexo II: Código de aplicación web}\label{anexo2}

\subsection*{Funciones}

\textbf{Caja XX}. Definición de \texttt{funcion}.

\lstinputlisting[language=R]{codigo/cabecera.R}
		
\end{document}